\documentclass[11pt,a4paper]{article}

% --- Packages ---
\usepackage[utf8]{inputenc}
\usepackage[T1]{fontenc}
\usepackage{lmodern}
\usepackage{microtype}
\usepackage[margin=1.1in, top=1.3in, bottom=1.3in]{geometry}
\usepackage[colorlinks=true, linkcolor=black, urlcolor=blue, citecolor=black]{hyperref}
\usepackage{booktabs}
\usepackage{tabularx}
\usepackage{longtable}
\usepackage{graphicx}
\usepackage{listings}
\usepackage{xcolor}
\usepackage{enumitem}
\usepackage{titlesec}
\usepackage{fancyhdr}
\usepackage{parskip}
\usepackage{array}
\usepackage{amsmath}

% --- Page style ---
\pagestyle{fancy}
\fancyhf{}
\lhead{\small Aroha Whitepaper v1.0}
\rhead{\small April 28, 2026 --- DRAFT}
\cfoot{\thepage}
\setlength{\headheight}{14pt}

% --- Section formatting ---
\titleformat{\section}{\large\bfseries}{\thesection.}{0.6em}{}[\vspace{2pt}\hrule\vspace{4pt}]
\titleformat{\subsection}{\normalsize\bfseries}{\thesubsection}{0.6em}{}
\titleformat{\subsubsection}{\normalsize\itshape}{\thesubsubsection}{0.6em}{}

% --- Code listings ---
\definecolor{codebg}{RGB}{248,248,248}
\definecolor{codecomment}{RGB}{100,100,100}
\definecolor{codekw}{RGB}{0,0,180}
\definecolor{codestr}{RGB}{160,20,20}
\definecolor{rulecolor}{RGB}{200,200,200}

\lstset{
  backgroundcolor=\color{codebg},
  basicstyle=\ttfamily\small,
  breaklines=true,
  breakatwhitespace=true,
  captionpos=b,
  commentstyle=\color{codecomment}\itshape,
  keywordstyle=\color{codekw}\bfseries,
  stringstyle=\color{codestr},
  frame=single,
  rulecolor=\color{rulecolor},
  showstringspaces=false,
  columns=flexible,
  keepspaces=true,
  tabsize=2,
  xleftmargin=12pt,
  xrightmargin=8pt,
  aboveskip=10pt,
  belowskip=6pt
}

% --- Title ---
\title{
  \vspace{-1cm}
  \textbf{\LARGE Aroha: Aroha Protocol} \\[0.6em]
  \large Whitepaper --- Specification v1.0 --- Draft
}
\author{
  Aroha Working Group \\
  \small Draft --- not submitted to any standards body
}
\date{April 28, 2026}

% ================================================================================
\begin{document}
% ================================================================================

\maketitle
\thispagestyle{fancy}

\begin{abstract}
Aroha (Aroha Protocol) is an open, decentralized, vendor-neutral
protocol for agent-to-agent communication, discovery, and atomic multi-agent
transactions. It defines how AI agents from different organizations discover each
other, establish cryptographic identity and trust, exchange messages securely, and
coordinate complex multi-step tasks that either fully succeed or fully roll back.
Aroha is language-agnostic and framework-agnostic. Any agent that correctly implements
these rules can interoperate with any other Aroha agent --- across vendors, clouds,
and programming languages. The protocol includes a deterministic policy engine for
prompt-injection defence, a semantic cache layer for inference cost reduction,
W3C Trace Context distributed tracing, and a no-flag-day migration path to NIST
post-quantum cryptography standards.
\end{abstract}

\tableofcontents
\newpage

% ================================================================================
\section{Introduction and Motivation}
% ================================================================================

\subsection{The Emerging Agent Economy}

AI agents are being deployed everywhere. Every enterprise, every startup, and every
developer team is building them. But today, every agent is an island. An agent built
by Salesforce cannot natively talk to an agent built by SAP. An agent running on
Anthropic's SDK cannot atomically transact with one running on OpenAI's. This is
the same structural problem the internet had before HTTP, that email had before SMTP,
and that payments had before ISO 8583.

\textbf{The agent economy is being built right now. If the interoperability problem
is not solved at the protocol layer, it will be solved by monopoly.}

\subsection{The Three Infrastructure Failures}

Building cross-vendor agent systems today results in three systematic failures:

\begin{table}[ht]
  \centering
  \renewcommand{\arraystretch}{1.6}
  \begin{tabularx}{\textwidth}{@{} >{\bfseries}p{2.5cm} p{5.0cm} X @{}}
    \toprule
    Failure & Phenomenon & Systemic Impact \\
    \midrule
    1. Discovery
      & No standard way for an agent to advertise what it can do, or for another
        agent to find it. Every connection is hardcoded by a human.
      & Multi-agent systems become brittle point-to-point integrations that do
        not scale. Humans become the bottleneck for agent coordination. \\
    2. Verification
      & No cryptographic proof of identity. Any message can be forged, replayed,
        or tampered with. Standard APIs trust callers implicitly via shared API keys.
      & Vulnerable to impersonation, replay, and man-in-the-middle attacks.
        Enterprises cannot safely deploy agents that communicate externally. \\
    3. Transactions
      & No standard rollback. When one step of a multi-agent saga fails,
        there is no compensation mechanism.
      & Multi-provider workflows leave partial state, duplicate charges, and
        disputes with no audit trail. \\
    \bottomrule
  \end{tabularx}
  \caption{The three dimensions of multi-agent infrastructure failure}
\end{table}

\subsection{Why Existing Solutions Fall Short}

\begin{table}[ht]
  \centering
  \renewcommand{\arraystretch}{1.4}
  \begin{tabularx}{\textwidth}{@{} l p{4.5cm} X @{}}
    \toprule
    \textbf{Solution} & \textbf{What it solves} & \textbf{What it misses} \\
    \midrule
    REST / OpenAPI      & HTTP call description       & No agent identity, no discovery, no transactions \\
    OAuth 2.0           & Delegated access for humans & Not designed for autonomous multi-agent flows \\
    MCP (Anthropic)     & LLM tool calling            & No agent-to-agent trust, no discovery, no saga \\
    A2A (Google)        & Agent-to-agent messaging    & No commerce layer, no capability-level reputation \\
    LangGraph / CrewAI  & In-framework orchestration  & Framework-local; no cross-vendor interoperability \\
    Blockchain alone    & Immutable records           & No application-layer protocol on top \\
    \bottomrule
  \end{tabularx}
  \caption{Gaps in existing agent communication solutions}
\end{table}

Aroha covers all four dimensions simultaneously:
\textbf{discovery + verification + transactions + commerce}.

% ================================================================================
\section{Design Principles}
% ================================================================================

Aroha is governed by six non-negotiable design principles:

\begin{enumerate}[leftmargin=*, label=\textbf{\arabic*.}]
  \item \textbf{Decentralized} --- No central authority controls agent discovery or
        registration. Any operator can run a conformant registry.
  \item \textbf{Self-sovereign identity} --- Every agent owns its cryptographic
        identity; no third party can revoke it.
  \item \textbf{Composable} --- Complex tasks are assembled from atomic,
        single-purpose agents that each do one thing well.
  \item \textbf{Atomic} --- Multi-agent transactions either fully succeed or fully
        roll back (saga pattern). Partial execution is structurally impossible.
  \item \textbf{Open standards} --- Built on W3C DIDs, JSON-LD, Ed25519, and
        optionally EVM smart contracts. No proprietary components.
  \item \textbf{Protocol-application separation} --- The protocol defines envelopes
        and handshakes. Applications decide what to do with messages. The kernel
        never imports extension packages.
\end{enumerate}

% ================================================================================
\section{Protocol Architecture}
% ================================================================================

\subsection{The Six-Layer Stack}

Aroha is a six-layer protocol stack. The core kernel (Layers 0--5) covers exactly
two things: \textit{agent discoverability} and \textit{agent verifiability}.
Everything else is a pluggable extension injected via middleware.

\begin{verbatim}
+------------------------------------------------------------------+
|  PROTOCOL EXTENSIONS  (never imported by kernel)                 |
|                                                                  |
|  @aroha-sdk/credentials    RBAC middleware, human credential issuance |
|  @aroha-sdk/settlement     NullSettlement | ApiKey | Quota | Escrow   |
|  @aroha-sdk/commerce       B2B agreements, B2C consent, C2B feedback  |
+------------------------------------------------------------------+
                          ^ mount via middleware / inject
+--------------------------------------------------+
|  Layer 5  Preference Context                     |
|           budget, consent tokens, constraints    |
+--------------------------------------------------+
|  Layer 4  Orchestration                          |
|           saga state machine: reserve/commit/cancel              |
+--------------------------------------------------+
|  Layer 3  Messaging                              |
|           signed JSON-LD envelopes, 15 types     |
+--------------------------------------------------+
|  Layer 2  Registry                               |
|           pluggable: HTTP | in-memory | EVM      |
|           + IPFS capability manifests            |
+--------------------------------------------------+
|  Layer 1  Identity and Auth                      |
|           did:aroha DIDs, Ed25519, W3C VCs        |
+--------------------------------------------------+
|  Layer 0  Transport                              |
|           HTTPS, WebSocket (wss://), TLS 1.3+    |
+--------------------------------------------------+
\end{verbatim}

\subsection{Agent Roles}

Aroha distinguishes two participant roles. A single agent may occupy both
simultaneously.

\begin{table}[ht]
  \centering
  \renewcommand{\arraystretch}{1.4}
  \begin{tabularx}{\textwidth}{@{} l X @{}}
    \toprule
    \textbf{Role} & \textbf{Responsibilities} \\
    \midrule
    \textbf{Orchestrator}
      & Receives user intents. Discovers providers via the registry. Drives the
        saga (sends Reserve, Commit, Cancel). Handles rollback on failure. \\
    \textbf{Provider}
      & Exposes one or more capabilities in a signed manifest. Participates in
        sagas by holding reservations and executing commits. \\
    \bottomrule
  \end{tabularx}
  \caption{Aroha agent roles}
\end{table}

\subsubsection{Message Flow Example}

\begin{verbatim}
User / UI
    |
    | POST /task  { "goal": "book NYC to Tokyo + 3 nights hotel" }
    v
Personal Agent (Orchestrator)
    |         |
    |         | ArohaRequest (signed envelope)
    |         +----> Flight Agent (Provider)
    |                   <---- ArohaResponse { flights: [...] }
    |
    | ArohaReserve { capabilityId: "hold-seat", params: {...} }
    +----> Hotel Agent (Provider)
    |         <---- ArohaReserveAck { reservationToken: "tok_abc" }
    |
    | ArohaCommit { reservationToken: "tok_abc" }
    +----> Hotel Agent
              <---- ArohaCommitAck { result: { confirmationId: "H-123" } }
\end{verbatim}

% ================================================================================
\section{Layer 0 --- Transport}
% ================================================================================

Every Aroha agent \textbf{MUST} expose three HTTP endpoints:

\begin{table}[ht]
  \centering
  \renewcommand{\arraystretch}{1.3}
  \begin{tabularx}{\textwidth}{@{} l l X @{}}
    \toprule
    \textbf{Method} & \textbf{Path} & \textbf{Purpose} \\
    \midrule
    \texttt{GET}  & \texttt{/.well-known/aroha-agent.json} & DID Document (public identity) \\
    \texttt{POST} & \texttt{/aroha/v1}                     & Inbound signed message endpoint \\
    \texttt{WS}   & \texttt{/aroha/stream?cid=<id>}        & Async streaming (ArohaStream) \\
    \bottomrule
  \end{tabularx}
  \caption{Required HTTP endpoints}
\end{table}

\textbf{Transport security requirements:}
\begin{itemize}[noitemsep]
  \item All connections \textbf{MUST} use TLS 1.3 or higher.
  \item Provider agents \textbf{SHOULD} use mutual TLS (mTLS).
  \item Plain HTTP \textbf{MUST NOT} be used in production.
  \item Localhost HTTP is permitted for development and automated testing only.
\end{itemize}

% ================================================================================
\section{Layer 1 --- Identity (did:aroha)}
% ================================================================================

\subsection{DID Method}

Every agent \textbf{MUST} have a W3C Decentralised Identifier:
\begin{center}
  \texttt{did:aroha:<unique-id>}
\end{center}

where \texttt{<unique-id>} matches \texttt{[a-z0-9][a-z0-9-]\{1,62\}}.
Examples: \texttt{did:aroha:flight-booking-agent}, \texttt{did:aroha:personal-agent-alice}.

\subsection{DID Document}

Served at \texttt{GET /.well-known/aroha-agent.json} and signed by the agent's
Ed25519 private key:

\begin{lstlisting}[caption=DID Document structure]
{
  "@context": [
    "https://www.w3.org/ns/did/v1",
    "https://w3id.org/security/suites/ed25519-2020/v1"
  ],
  "id": "did:aroha:my-agent",
  "verificationMethod": [{
    "id": "did:aroha:my-agent#key-1",
    "type": "Ed25519VerificationKey2020",
    "controller": "did:aroha:my-agent",
    "publicKeyMultibase": "z<base58btc-encoded-32-bytes>"
  }],
  "authentication":  ["did:aroha:my-agent#key-1"],
  "assertionMethod": ["did:aroha:my-agent#key-1"],
  "service": [{
    "id": "did:aroha:my-agent#aroha",
    "type": "ArohaEndpoint",
    "serviceEndpoint": "https://my-agent.example.com/aroha/v1"
  }],
  "created": "2026-04-28T00:00:00Z"
}
\end{lstlisting}

\subsection{Trust Levels}

\begin{table}[ht]
  \centering
  \renewcommand{\arraystretch}{1.3}
  \begin{tabularx}{\textwidth}{@{} >{\bfseries}l p{5.5cm} X @{}}
    \toprule
    Level & How to achieve & Permitted operations \\
    \midrule
    0 & None                                        & Read manifests, query registry \\
    1 & On-chain registration + minimum stake       & Search, negotiate \\
    2 & W3C VC (CorporateIdentityCredential)        & Reserve (low-value) \\
    3 & On-chain reputation score >= 8000 / 10000   & Reserve + commit (high-value) \\
    4 & CapabilityAuthorizationCredential           & Full access + payments \\
    \bottomrule
  \end{tabularx}
  \caption{Agent trust levels and permitted operations}
\end{table}

\subsection{Verifiable Credentials}

Agents \textbf{MAY} carry W3C Verifiable Credentials to prove claims beyond their DID:

\begin{table}[ht]
  \centering
  \renewcommand{\arraystretch}{1.3}
  \begin{tabularx}{\textwidth}{@{} l l X @{}}
    \toprule
    \textbf{VC Type} & \textbf{Required Trust Level} & \textbf{Purpose} \\
    \midrule
    \texttt{CorporateIdentityCredential}       & 2 & Proves corporate operator \\
    \texttt{ComplianceCertificationCredential} & 2 & Proves regulatory compliance \\
    \texttt{CapabilityAuthorizationCredential} & 4 & Authorizes capabilities and payments \\
    \bottomrule
  \end{tabularx}
  \caption{W3C Verifiable Credential types used in Aroha}
\end{table}

% ================================================================================
\section{Layer 2 --- Registry}
% ================================================================================

\subsection{Registry Backends}

Any conformant registry implements \texttt{IArohaRegistry}. Four backends are
defined and are interchangeable at the SDK level:

\begin{table}[ht]
  \centering
  \renewcommand{\arraystretch}{1.4}
  \begin{tabularx}{\textwidth}{@{} l l X @{}}
    \toprule
    \textbf{Implementation} & \textbf{Backend} & \textbf{Use case} \\
    \midrule
    Managed HTTP Registry     & REST API (SaaS or self-hosted) & Enterprise default. No blockchain required. \\
    Self-hosted HTTP Registry & REST API (operator-run)        & Regulated industries (healthcare, finance). \\
    In-Memory Registry        & In-process Maps                & Local dev, tests, intra-enterprise networks. \\
    Blockchain Registry       & EVM smart contract             & Trustless, permissionless open ecosystems. \\
    \bottomrule
  \end{tabularx}
  \caption{Pluggable registry implementations}
\end{table}

\subsection{HTTP Registry REST Endpoints}

\begin{table}[ht]
  \centering
  \renewcommand{\arraystretch}{1.3}
  \begin{tabularx}{\textwidth}{@{} l l X @{}}
    \toprule
    \textbf{Method} & \textbf{Path} & \textbf{Purpose} \\
    \midrule
    \texttt{POST}   & \texttt{/v1/agents}                   & Register an agent \\
    \texttt{DELETE} & \texttt{/v1/agents/:didHash}          & Deactivate an agent \\
    \texttt{PUT}    & \texttt{/v1/agents/:didHash/manifest} & Update capability manifest CID \\
    \texttt{GET}    & \texttt{/v1/agents/:didHash}          & Resolve agent record by DID \\
    \texttt{GET}    & \texttt{/v1/agents?capabilities=a,b}  & Query agents by capability IDs \\
    \bottomrule
  \end{tabularx}
  \caption{HTTP registry REST endpoints}
\end{table}

\subsection{Capability Manifest}

Published to IPFS; CID anchored in the registry. The manifest is Ed25519-signed.
Any manifest with an invalid signature is silently rejected by \texttt{ArohaDiscovery}:

\begin{lstlisting}[caption=Capability manifest format]
{
  "aroha": "1.0",
  "did": "did:aroha:flight-agent",
  "name": "FlightBookingAgent",
  "version": "2.1.0",
  "endpoint": "https://flights.example.com/aroha/v1",
  "capabilities": [
    {
      "id": "book-flight",
      "version": "1.0.0",
      "async": false,
      "maxResponseMs": 5000,
      "requiredTrustLevel": 3,
      "pricing": {
        "model": "per-transaction",
        "amount": 0.05,
        "currency": "USD"
      },
      "sagaRole": "participant",
      "reserveCapabilityId": "hold-seat",
      "cancelCapabilityId": "release-seat"
    }
  ],
  "protocols": ["aroha/1.0"],
  "signature": "<base64url-Ed25519-signature>"
}
\end{lstlisting}

\subsection{Reputation Oracle}

Reputation scores (0--10000) are maintained by the \texttt{ArohaReputation} oracle.
The exponential moving average formula ensures recent behavior dominates:

\[
  \text{newScore} = \left\lfloor \text{oldScore} \times 0.85 + \text{signal} \times 0.15 \right\rfloor
\]

\begin{table}[ht]
  \centering
  \renewcommand{\arraystretch}{1.3}
  \begin{tabularx}{\textwidth}{@{} l l X @{}}
    \toprule
    \textbf{Outcome} & \textbf{Signal value} & \textbf{Notes} \\
    \midrule
    \texttt{SUCCESS}        & 10000 & Committed and result delivered \\
    \texttt{COMPENSATED}    &  3000 & Saga rolled back cleanly \\
    \texttt{FAILURE}        &  1000 & Timed out or returned error \\
    \texttt{DISPUTE\_LOST}  &     0 & Provider found at fault \\
    \bottomrule
  \end{tabularx}
  \caption{Reputation signal values. New agents start at 5000 (neutral).}
\end{table}

% ================================================================================
\section{Layer 3 --- Messaging}
% ================================================================================

\subsection{Envelope Structure}

Every Aroha message uses the same outer signed envelope:

\begin{lstlisting}[caption=Standard Aroha signed envelope]
{
  "@context": [
    "https://aroha-labs.com/contexts/v1",
    "https://w3id.org/security/suites/ed25519-2020/v1"
  ],
  "id":            "urn:uuid:<uuid-v4>",
  "type":          "<MessageType>",
  "from":          "did:aroha:<sender>",
  "to":            "did:aroha:<recipient>",
  "created":       "2026-04-28T08:00:00Z",
  "expires":       "2026-04-28T08:05:00Z",
  "nonce":         "<cryptographically-random-32-byte-hex>",
  "correlationId": "<string>",
  "body": { },
  "proof": {
    "type":               "Ed25519Signature2020",
    "created":            "2026-04-28T08:00:00Z",
    "verificationMethod": "did:aroha:<sender>#key-1",
    "proofPurpose":       "authentication",
    "proofValue":         "<base64url-Ed25519-signature>"
  }
}
\end{lstlisting}

\texttt{expires} defaults to \texttt{created + 300s}. Nonces are validated against
a rolling 5-minute window; replay is structurally impossible.

\subsection{The 15 Message Types}

Aroha defines exactly 15 message types. Applications \textbf{MUST NOT} invent new
top-level types --- extension is done through the \texttt{body} field only.

\begin{longtable}{@{} l l c X @{}}
  \toprule
  \textbf{Type} & \textbf{Direction} & \textbf{Saga} & \textbf{Description} \\
  \midrule
  \endhead
  \texttt{ArohaRequest}            & Orch $\to$ Prov  &   & Invoke a read-only capability \\
  \texttt{ArohaResponse}           & Prov $\to$ Orch  &   & Capability result \\
  \texttt{ArohaStream}             & Prov $\to$ Orch  &   & Async streaming (WebSocket) \\
  \texttt{ArohaNegotiate}          & Orch $\to$ Prov  &   & Open a price negotiation \\
  \texttt{ArohaCounterOffer}       & Prov $\to$ Orch  &   & Provider counter-proposes price \\
  \texttt{ArohaAccept}             & Orch $\to$ Prov  &   & Accept current offer \\
  \texttt{ArohaReserve}            & Orch $\to$ Prov  & Y & Hold inventory \\
  \texttt{ArohaReserveAck}         & Prov $\to$ Orch  & Y & Reservation token confirmed \\
  \texttt{ArohaCommit}             & Orch $\to$ Prov  & Y & Finalize reservation (idempotent) \\
  \texttt{ArohaCommitAck}          & Prov $\to$ Orch  & Y & Commitment confirmed + result \\
  \texttt{ArohaCancel}             & Orch $\to$ Prov  & Y & Compensating transaction / rollback \\
  \texttt{ArohaCancelAck}          & Prov $\to$ Orch  & Y & Cancellation confirmed \\
  \texttt{ArohaDelegate}           & Orch $\to$ Orch  &   & Forward to sub-orchestrator \\
  \texttt{ArohaError}              & Prov $\to$ Orch  &   & Typed, retryable error \\
  \texttt{ArohaSatisfactionSignal} & Orch $\to$ Prov  &   & Post-saga reputation feedback \\
  \bottomrule
  \caption{Complete Aroha message type reference}
\end{longtable}

% ================================================================================
\section{Layer 4 --- Orchestration and the Saga Pattern}
% ================================================================================

\subsection{The Saga State Machine}

The saga pattern is a Aroha protocol primitive. Every multi-step transaction passes
through the following finite-state machine:

\begin{verbatim}
                          ArohaReserve sent
[*] --> [IDLE] --------------------------> [RESERVING]
                                                |
                                 ArohaReserveAck +---> timeout / ArohaError
                                    |                        |
                               [RESERVED]              [FAILED] --> [*]
                                    |
                      +-------------+------------------+
                      |                                |
              ArohaCommit sent                  orchestrator aborts
                      |                                |
                [COMMITTING]              [COMPENSATING] <--.
                      |                       |              |
          ArohaCommitAck                ArohaCancelAck         |
                  |                          |          commit fails
           [COMMITTED]               [COMPENSATED]           |
               [*]                       [*]                 |
                            +--------------------------------+
\end{verbatim}

\subsection{Execution Phases}

\begin{enumerate}[label=\textbf{Phase \arabic*:}]
  \item \textbf{RESERVE (parallel)} --- The orchestrator sends \texttt{ArohaReserve}
        to all participants simultaneously. Each provider holds the requested
        resource for \texttt{holdDurationMs} milliseconds. Any failure aborts and
        triggers LIFO compensation immediately.

  \item \textbf{AWAITING APPROVAL (optional)} --- If a \texttt{HumanApprovalGate}
        is configured, the saga pauses after all reserves succeed. A human decision
        is polled via webhook. Rejection triggers LIFO compensation; no money moves.

  \item \textbf{COMMIT (sequential)} --- The orchestrator sends \texttt{ArohaCommit}
        to each participant in order. Providers \textbf{MUST} implement commit
        idempotently using \texttt{reservationToken} as the database key.

  \item \textbf{COMPENSATE (LIFO, on failure)} --- If any commit fails, the
        orchestrator sends \texttt{ArohaCancel} to all previously committed
        participants in reverse order. Cancel handlers \textbf{MUST} always succeed.
\end{enumerate}

\subsection{Read-Only Steps}

Steps marked \texttt{readOnly: true} use \texttt{ArohaRequest} and are skipped
entirely during commit and compensation. Useful for enrichment queries, sentiment
lookups, and price checks that should not hold reservations.

\subsection{Human-in-the-Loop Approval}

The \texttt{WebhookApprovalGate} pauses the saga after all reserves succeed, posts
a signed JSON payload to any webhook (Slack, Teams, or custom UI), and polls for
a decision. Rejection triggers full LIFO compensation instantly.

% ================================================================================
\section{Layer 5 --- Preference Context}
% ================================================================================

Every \texttt{ArohaRequest} and \texttt{ArohaReserve} \textbf{MAY} carry a
\texttt{preferenceContext} block encoding user constraints into the protocol payload:

\begin{lstlisting}[caption=Preference context block inside a ArohaRequest or ArohaReserve body]
"preferenceContext": {
  "budget":       { "max": 500, "currency": "USD"         },
  "quality":      { "minRating": 4.0                      },
  "constraints":  [ "non-stop-only", "window-seat"        ],
  "consentToken": "<signed-GDPR-consent-token>"
}
\end{lstlisting}

Providers \textbf{MUST} respect \texttt{budget.max}. A provider that returns a
result exceeding the declared budget \textbf{MUST} return
\texttt{ArohaError\{code: Aroha\_BUDGET\_EXCEEDED\}}.

% ================================================================================
\section{Security Model}
% ================================================================================

Aroha is designed for zero-trust agent communication. Every claim is
cryptographically verifiable before any application code runs.

\subsection{Threat Model}

\begin{table}[ht]
  \centering
  \renewcommand{\arraystretch}{1.5}
  \begin{tabularx}{\textwidth}{@{} l X @{}}
    \toprule
    \textbf{Threat} & \textbf{Mitigation} \\
    \midrule
    Message impersonation     & Ed25519 signature on every envelope \\
    Replay attacks            & 32-byte random nonce + \texttt{expires} TTL; rolling 5-min store \\
    Message tampering         & Signature covers canonicalized body (sorted keys, proof omitted) \\
    Wrong recipient           & \texttt{to} DID validated before \texttt{onMessage} fires \\
    MITM / eavesdropping      & TLS 1.3+ mandatory on all connections \\
    Privilege escalation      & RBAC enforced before application code \\
    Credential forgery        & Credentials signed by issuer's Ed25519 key \\
    Stale credentials         & \texttt{expiresAt} enforced; registry credentials are revocable \\
    Registry manipulation     & Manifest signature verified against public key \\
    Partial saga / dirty state& LIFO compensation guarantees no partial commits \\
    \bottomrule
  \end{tabularx}
  \caption{Zero-trust threat model and mitigations}
\end{table}

\subsection{Message Signing Algorithm}

\begin{lstlisting}[caption=Ed25519 signing and verification]
// --- Signing (sender side) ---
signingInput = UTF8(
  JSON.stringify(
    sortedKeys(message, omit="proof")  // canonical form
  )
)
proofValue = base64url(
  Ed25519Sign(signingInput, agentPrivateKey)
)

// --- Verification (receiver side) ---
isValid = await ed.verifyAsync(
  fromBase64url(envelope.proof.proofValue), // signature bytes
  new TextEncoder().encode(                  // message bytes
    JSON.stringify(sortKeys(bodyWithoutProof))
  ),
  senderPublicKey  // resolved from registry or DID Document
)
\end{lstlisting}

\subsection{RBAC Enforcement Order}

When \texttt{@aroha-sdk/credentials} is mounted, the server evaluates in this order
before \texttt{onMessage} is called:

\begin{enumerate}[noitemsep]
  \item If \texttt{body.credentialId} present: resolve from \texttt{CredentialRegistry},
        verify signature, extract roles.
  \item If \texttt{body.credentialToken} present: verify inline credential signature,
        extract roles.
  \item Else: resolve from \texttt{rbacPolicy} options, use agent roles.
\end{enumerate}

\subsection{Role--Action Permission Matrix}

\begin{table}[ht]
  \centering
  \renewcommand{\arraystretch}{1.3}
  \begin{tabularx}{\textwidth}{@{} l X l @{}}
    \toprule
    \textbf{Role} & \textbf{Permitted actions} & \textbf{Notes} \\
    \midrule
    \texttt{human:admin}        & ALL                                            & Full access \\
    \texttt{human:user}         & query, reserve, commit, cancel, negotiate      & Standard user \\
    \texttt{human:readonly}     & query                                          & Read-only \\
    \texttt{agent:orchestrator} & query, reserve, commit, cancel, delegate, negotiate & Cross-agent \\
    \texttt{agent:provider}     & receive only                                   & No restriction \\
    \texttt{agent:observer}     & query                                          & Monitoring \\
    \bottomrule
  \end{tabularx}
  \caption{RBAC role-action permission matrix}
\end{table}

\subsection{Production Security Checklist}

\begin{itemize}[noitemsep]
  \item[$\square$] TLS 1.3 on all endpoints
  \item[$\square$] mTLS between known peer agents
  \item[$\square$] Private key in HSM or secrets manager, never in env vars or source
  \item[$\square$] Credential TTLs $\leq$ 1 hour for \texttt{human:user}
  \item[$\square$] Registry-based credentials for all production sessions (enables revocation)
  \item[$\square$] Rate limiter in front of \texttt{ArohaServer}
  \item[$\square$] Monitor for \texttt{Aroha\_UNAUTHORIZED} spikes
\end{itemize}

% ================================================================================
\section{Extension: Settlement (@aroha-sdk/settlement)}
% ================================================================================

The settlement layer is entirely optional. The protocol kernel has no opinion on
whether agent interactions are free, metered, or settled via smart contracts.
Settlement strategies hook into the saga via \texttt{onReserve}, \texttt{onCommit},
\texttt{onCancel}, and \texttt{onSagaComplete}.

\subsection{Settlement Strategies}

\begin{table}[ht]
  \centering
  \renewcommand{\arraystretch}{1.4}
  \begin{tabularx}{\textwidth}{@{} l X X @{}}
    \toprule
    \textbf{Strategy} & \textbf{Mechanism} & \textbf{Best for} \\
    \midrule
    \texttt{NullSettlement}     & No-op. All operations free.
      & Internal networks, development \\
    \texttt{ApiKeySettlement}   & Validates API keys per provider.
      & SaaS provider integration \\
    \texttt{QuotaSettlement}    & Holds quota on Reserve; deducts on Commit; releases on Cancel.
      & Rate-limited enterprise APIs \\
    \texttt{EscrowSettlement}   & Opens on-chain escrow on Reserve; approves on Commit; refunds on Cancel.
      & Trustless cross-org payments \\
    \texttt{PostgresSettlement} & Tracks saga lifecycle in Postgres.
      & Enterprise Web2 deployments \\
    \texttt{StripeSettlement}   & Stripe manual-capture PaymentIntent per saga step.
      & SaaS consumer billing \\
    \bottomrule
  \end{tabularx}
  \caption{Built-in settlement strategies}
\end{table}

\subsection{Escrow Settlement --- Invisible Web3}

\begin{enumerate}[noitemsep]
  \item Developer calls \texttt{wallet.buyData(providerDID, \{budget, capability\})}
  \item \texttt{WalletService} generates ephemeral session key; Paymaster sponsors gas
  \item \texttt{ArohaReserve} sent; \texttt{onReserve} opens escrow with funds held
  \item \texttt{ArohaCommit} sent; \texttt{onCommit} approves split for provider
  \item \texttt{onSagaComplete(SUCCESS)} releases funds to provider
  \item On failure: \texttt{onSagaComplete(COMPENSATED)} refunds orchestrator in full
\end{enumerate}

\subsection{Stripe Settlement Lifecycle}

\begin{table}[ht]
  \centering
  \renewcommand{\arraystretch}{1.3}
  \begin{tabularx}{\textwidth}{@{} l l X @{}}
    \toprule
    \textbf{Saga phase} & \textbf{Stripe action} & \textbf{Effect} \\
    \midrule
    \texttt{ArohaReserve} & PaymentIntent created (\texttt{capture\_method: "manual"}) & Authorization only; no charge \\
    \texttt{ArohaCommit}  & PaymentIntent captured                                     & Money moves \\
    \texttt{ArohaCancel}  & PaymentIntent cancelled                                    & Hold released; zero charge \\
    \bottomrule
  \end{tabularx}
  \caption{Stripe settlement lifecycle}
\end{table}

% ================================================================================
\section{Extension: Commerce (@aroha-sdk/commerce)}
% ================================================================================

\begin{table}[ht]
  \centering
  \renewcommand{\arraystretch}{1.4}
  \begin{tabularx}{\textwidth}{@{} l l X @{}}
    \toprule
    \textbf{Pattern} & \textbf{Relationship} & \textbf{Example} \\
    \midrule
    B2B & Org $\leftrightarrow$ Org
      & Logistics agent consumes freight-pricing agent under quarterly contract + SLA \\
    B2C & Org $\to$ Consumer
      & Travel agent serves personal agent under monthly subscription + GDPR consent \\
    C2B & Consumer $\to$ Org
      & Personal agent submits quality feedback; recorded on-chain for reputation \\
    \bottomrule
  \end{tabularx}
  \caption{Commerce relationship patterns}
\end{table}

\subsection{B2B Agreement Lifecycle}

\begin{enumerate}[noitemsep, label=\textbf{\arabic*.}]
  \item Provider creates draft with capabilities, billing terms, and SLA terms
  \item Provider signs with Ed25519 private key
  \item Consumer countersigns --- agreement becomes active
  \item Both parties can \texttt{verifyAgreement()} at any time
  \item \texttt{B2BBillingSettlement} enforces quotas and SLA per call
  \item \texttt{generateInvoice()} produces a line-item invoice from billing records
\end{enumerate}

\subsection{B2C Consent (GDPR-aligned)}

Every consent record is signed by the consumer's Ed25519 key. Providers
\textbf{MUST} verify the signature before using data for any declared purpose.
Allowed use purposes, retention days, cross-border transfer flags, and GDPR
applicability are encoded in the signed record.

% ================================================================================
\section{Extension: Credentials and RBAC (@aroha-sdk/credentials)}
% ================================================================================

Human users authenticate out-of-band (password, OAuth, passkey) and receive a
signed \texttt{HumanCredential} with a short TTL. Two delivery modes:

\begin{itemize}[noitemsep]
  \item \textbf{Inline token} (\texttt{credentialToken}): attached to the message
        body. Cannot be revoked before \texttt{expiresAt}. Use TTL $\leq$ 1 hour.
  \item \textbf{Registry reference} (\texttt{credentialId}): only the UUID is sent;
        the agent resolves it from the \texttt{CredentialRegistry} API. Supports
        instant revocation via \texttt{DELETE /aroha/credentials/:id}.
\end{itemize}

\begin{table}[ht]
  \centering
  \renewcommand{\arraystretch}{1.3}
  \begin{tabularx}{\textwidth}{@{} l l X l @{}}
    \toprule
    \textbf{Method} & \textbf{Path} & \textbf{Purpose} & \textbf{Auth} \\
    \midrule
    \texttt{POST}   & \texttt{/aroha/credentials}         & Register credential & Signature \\
    \texttt{GET}    & \texttt{/aroha/credentials/:id}     & Resolve by ID       & Public \\
    \texttt{GET}    & \texttt{/aroha/credentials?userId=} & List by user        & Public \\
    \texttt{DELETE} & \texttt{/aroha/credentials/:id}     & Revoke              & Revoke proof \\
    \bottomrule
  \end{tabularx}
  \caption{Credential Registry HTTP endpoints}
\end{table}

% ================================================================================
\section{Error Codes}
% ================================================================================

All Aroha errors are returned as a \texttt{ArohaError} envelope with a typed
\texttt{code}, a human-readable \texttt{message}, and a \texttt{retryable} boolean.
Callers \textbf{MUST NOT} retry a non-retryable error without fixing the root cause.

\begin{longtable}{@{} l c X @{}}
  \toprule
  \textbf{Code} & \textbf{Retryable} & \textbf{Meaning} \\
  \midrule
  \endhead
  \texttt{Aroha\_UNAUTHORIZED}         & no  & No credential supplied or signature verification failed \\
  \texttt{Aroha\_FORBIDDEN}            & no  & Credential valid but role lacks the required permission \\
  \texttt{Aroha\_CREDENTIAL\_EXPIRED}  & no  & \texttt{HumanCredential.expiresAt} is in the past \\
  \texttt{Aroha\_CREDENTIAL\_REVOKED}  & no  & Credential has been revoked via \texttt{DELETE /aroha/credentials/:id} \\
  \texttt{Aroha\_INVALID\_SIGNATURE}   & no  & Ed25519 proof verification failed \\
  \texttt{Aroha\_REPLAY\_DETECTED}     & no  & Nonce already seen within the rolling 5-minute window \\
  \texttt{Aroha\_MESSAGE\_EXPIRED}     & no  & \texttt{envelope.expires} is in the past \\
  \texttt{Aroha\_WRONG\_RECIPIENT}     & no  & \texttt{to} DID does not match this agent's DID \\
  \texttt{Aroha\_UNKNOWN\_CAPABILITY}  & no  & Capability ID is not supported by this provider \\
  \texttt{Aroha\_INVALID\_TOKEN}       & no  & Reservation token is unknown or has already been cancelled \\
  \texttt{Aroha\_RESERVATION\_EXPIRED} & no  & \texttt{holdDurationMs} elapsed before commit was received \\
  \texttt{Aroha\_NO\_INVENTORY}        & yes & Capability is temporarily unavailable; safe to retry later \\
  \texttt{Aroha\_ALREADY\_COMMITTED}   & no  & Commit received for a reservation that was already cancelled \\
  \texttt{Aroha\_UPSTREAM\_ERROR}      & yes & Provider failed to reach an upstream dependency \\
  \texttt{Aroha\_INTERNAL\_ERROR}      & yes & Unexpected internal error; check provider logs \\
  \texttt{Aroha\_RATE\_LIMITED}        & yes & Rate limit exceeded; respect \texttt{Retry-After} in \texttt{details} \\
  \bottomrule
  \caption{Complete Aroha error code reference}
\end{longtable}

% ================================================================================
\section{SDK Package Reference}
% ================================================================================

The Aroha SDK is a TypeScript-first monorepo. Every package is independently
installable via npm. The kernel packages (\texttt{@aroha-sdk/core} through
\texttt{@aroha-sdk/preferences}) have zero runtime dependencies on extension packages.

\begin{table}[ht]
  \centering
  \renewcommand{\arraystretch}{1.4}
  \begin{tabularx}{\textwidth}{@{} l l X @{}}
    \toprule
    \textbf{Package} & \textbf{Layer} & \textbf{Provides} \\
    \midrule
    \texttt{@aroha-sdk/core}             & 0--3 & Transport, DID identity, Ed25519 crypto, signed messaging \\
    \texttt{@aroha-sdk/registry}         & 2    & HTTP, in-memory, and EVM registry clients; IPFS pinning; capability discovery \\
    \texttt{@aroha-sdk/orchestrator}     & 4    & \texttt{SagaEngine}, \texttt{WebhookApprovalGate}, read-only step support \\
    \texttt{@aroha-sdk/preferences}      & 5    & Preference context builders, validators, consent token issuance \\
    \texttt{@aroha-sdk/credentials}      & ext  & RBAC server middleware, \texttt{HumanCredential} issuance, credential registry \\
    \texttt{@aroha-sdk/settlement}       & ext  & \texttt{NullSettlement}, \texttt{ApiKeySettlement}, \texttt{QuotaSettlement}, \texttt{EscrowSettlement}, \texttt{PostgresSettlement}, \texttt{StripeSettlement} \\
    \texttt{@aroha-sdk/commerce}         & ext  & B2B signed agreements, B2C subscriptions and GDPR consent, C2B feedback, audit log \\
    \texttt{@aroha-sdk/telemetry}        & ext  & W3C Trace Context propagation, distributed saga tracing, OpenTelemetry integration \\
    \texttt{@aroha-sdk/policy}           & ext  & Deterministic JSON policy enforcement; prompt-injection surface closure \\
    \texttt{@aroha-sdk/cache}            & ext  & SHA-256 content-addressed semantic cache; in-memory and Redis backends \\
    \texttt{@aroha-sdk/trusted-mesh}     & ext  & Bypass Ed25519 verification for allowlisted DIDs inside a VPC or mTLS mesh \\
    \texttt{@aroha-sdk/mcp-bridge}       & ext  & Expose any MCP server's tools as typed Aroha capability providers \\
    \texttt{@aroha-sdk/a2a-bridge}       & ext  & Bidirectional translation between Google A2A messages and Aroha envelopes \\
    \texttt{@aroha-sdk/langchain-bridge} & ext  & Expose Aroha capabilities as LangChain tools; wrap LangChain agents as providers \\
    \bottomrule
  \end{tabularx}
  \caption{Aroha SDK package reference (14 packages)}
\end{table}

% ================================================================================
\section{Interoperability and Bridges}
% ================================================================================

Aroha is designed to complement, not replace, existing agent ecosystems. All bridges
present a standard Aroha interface outward, so orchestrators do not need to know which
underlying framework a provider uses.

\begin{itemize}[leftmargin=*, itemsep=4pt]
  \item \textbf{MCP Bridge} (\texttt{@aroha-sdk/mcp-bridge}): Wraps any Model Context
        Protocol server as a Aroha capability provider. MCP tool calls arrive wrapped
        in a standard signed Aroha envelope, enabling per-call micro-payments and
        capability-level reputation scoring for MCP data sources --- with zero changes
        to the MCP server itself.

  \item \textbf{A2A Bridge} (\texttt{@aroha-sdk/a2a-bridge}): Translates Google's
        Agent-to-Agent (A2A) task and message objects into Aroha signed envelopes and
        vice versa. An A2A agent can participate as a provider in a Aroha saga without
        modification.

  \item \textbf{LangChain Bridge} (\texttt{@aroha-sdk/langchain-bridge}): Exposes
        Aroha capabilities as \texttt{StructuredTool} instances for LangChain agents,
        and wraps LangChain \texttt{AgentExecutor} instances as Aroha providers. Minimal
        changes to existing LangChain code are required.

  \item \textbf{Trusted Mesh} (\texttt{@aroha-sdk/trusted-mesh}): For agents that
        communicate exclusively within a private Kubernetes cluster or VPC behind
        mutual TLS, the \texttt{TrustedMeshMiddleware} can skip Ed25519 signature
        verification for a configurable allowlist of DIDs. Nonce deduplication
        and message expiry checks are always enforced regardless of trust mode.
\end{itemize}

% ================================================================================
\section{Specification Revision History}
% ================================================================================

\begin{table}[ht]
  \centering
  \renewcommand{\arraystretch}{1.4}
  \begin{tabularx}{\textwidth}{@{} l X @{}}
    \toprule
    \textbf{Date} & \textbf{Changes} \\
    \midrule
    2026-04-14
      & Initial draft (not submitted to the W3C or any other body). Core kernel (Layers 0--5)
        specified. \texttt{did:aroha} method defined. 15 message types enumerated.
        Saga FSM formalized. \\
    2026-04-21
      & \texttt{@aroha-sdk/commerce} extension finalized (B2B agreements, B2C consent,
        C2B feedback). \texttt{@aroha-sdk/credentials} extracted from kernel into
        separate extension package. \texttt{IArohaRegistry} pluggable backend
        interface finalized with four implementations. \\
    2026-04-28
      & Kernel hardening: \texttt{auth/} directory removed from \texttt{@aroha-sdk/core};
        no extension package is imported by the kernel. \texttt{WalletService}
        relocated to \texttt{@aroha-sdk/settlement/payments}. Ed25519$\to$X25519 ECDH
        key conversion algorithm corrected. Full Vitest test suite added
        (100+ tests across all 10 packages). SDK documentation complete. \\
    \bottomrule
  \end{tabularx}
  \caption{Specification revision history}
\end{table}

% ================================================================================
\section{Extension: Policy Engine (@aroha-sdk/policy)}
% ================================================================================

\subsection{Motivation: The LLM Security Boundary Problem}

LLMs are probabilistic. A sufficiently crafted prompt can cause a model to
construct a \texttt{transfer-funds} request for an amount orders of magnitude
above the operator's policy limit, embed a prompt-injection string inside a
structured parameter to manipulate a downstream agent, or add an unexpected
\texttt{admin: true} field to a capability that only allows read access.

\texttt{@aroha-sdk/policy} closes this surface by enforcing deterministic, rule-based
constraints \textit{after} Ed25519 signature verification and \textit{before}
the \texttt{onMessage} handler dispatches to any LLM or application code.
No reasoning is involved in policy evaluation --- rules are pure JSON predicates.

\subsection{Policy Rule Types}

\begin{table}[ht]
  \centering
  \renewcommand{\arraystretch}{1.4}
  \begin{tabularx}{\textwidth}{@{} l X @{}}
    \toprule
    \textbf{Rule} & \textbf{Enforcement} \\
    \midrule
    \texttt{maxSpendPerCall}
      & Reject requests where a numeric parameter exceeds a configured limit.
        Returns \texttt{Aroha\_FORBIDDEN} with the limit and actual value. \\
    \texttt{allowedParamKeys}
      & Reject requests containing any parameter key not on an explicit whitelist.
        Prevents undocumented or injected fields from reaching application code. \\
    \texttt{blockedParamValues}
      & Reject requests containing specific key/value combinations
        (e.g.\ \texttt{action=delete}, \texttt{admin=true}). \\
    \texttt{requiredParamKeys}
      & Reject requests missing mandatory fields. Prevents incomplete
        or LLM-hallucinated partial requests from executing. \\
    \texttt{paramAllowedValues}
      & Allowlist enforcement for enum-style parameters
        (e.g.\ \texttt{env} must be \texttt{staging} or \texttt{production}). \\
    \texttt{rateLimitPerMinute}
      & Sliding-window rate limit per sender DID per capability. \\
    \texttt{requiresHITLApproval}
      & Marks a capability as always requiring human sign-off, independent
        of \texttt{SagaEngine} configuration. \\
    \bottomrule
  \end{tabularx}
  \caption{Policy rule types supported by \texttt{ArohaPolicyEngine}}
\end{table}

\subsection{Prompt-Injection Surface Closure}

The companion \texttt{createStructuredPayloadMiddleware} rejects messages
containing string parameters longer than a configured maximum (default 512
characters) or nested parameter objects deeper than a configured depth (default 3
levels). This structurally prevents prompt-injection strings from reaching any
downstream LLM call while imposing no burden on well-formed agent requests.

Policy evaluation order: \textbf{signature} $\to$ \textbf{RBAC}
$\to$ \textbf{policy engine} $\to$ \textbf{structured payload check}
$\to$ \textbf{onMessage}.

% ================================================================================
\section{Extension: Semantic Cache (@aroha-sdk/cache)}
% ================================================================================

\subsection{The Agent Token-Cost Problem}

Agent-to-agent networks amplify inference costs. A saga involving five agents,
each making two inter-agent calls, incurs ten times the inference cost of a
single-agent workflow. Repeated lookups --- open issues, product availability,
flight prices, reputation scores --- spend tokens on questions whose answers have
not changed since the previous call.

\texttt{@aroha-sdk/cache} eliminates redundant inference by caching
\texttt{ArohaRequest}$\to$\texttt{ArohaResponse} pairs with content-addressed
keys. The second through \textit{n}th calls for the same capability with the
same canonical parameters return instantly from cache, with zero LLM invocations.

\subsection{Caching Mechanics}

\begin{itemize}[noitemsep]
  \item \textbf{Cache key}: SHA-256 of
        \texttt{capabilityId + canonical JSON of params} (key-order independent).
        Identical requests from different orchestrators, with different nonces or
        correlation IDs, hash to the same key.
  \item \textbf{Client-side} (\texttt{CachingArohaClient}): Intercepts
        \texttt{send()} before the network. Cache hit returns a synthesized
        \texttt{ArohaResponse} immediately.
  \item \textbf{Server-side} (\texttt{withCaching}): Wraps \texttt{onMessage}.
        Cache hit builds and signs a fresh \texttt{ArohaResponse} from stored data.
  \item \textbf{Saga safety}: Only \texttt{ArohaRequest} messages are cached.
        \texttt{ArohaReserve}, \texttt{ArohaCommit}, and \texttt{ArohaCancel} always
        pass through --- they carry side effects that must not be replayed from cache.
  \item \textbf{Backends}: \texttt{InMemorySemanticCache} (development) and
        \texttt{RedisSemanticCache} (production). Both implement the same
        \texttt{ISemanticCache} interface.
  \item \textbf{TTL}: Configurable per capability. Example: 60 seconds for
        \texttt{search-flights}, 10 minutes for \texttt{get-issue}.
  \item \textbf{Invalidation}: \texttt{cache.invalidate(capabilityId)} after any
        write operation that would stale the cached result.
\end{itemize}

% ================================================================================
\section{Extension: Distributed Tracing (@aroha-sdk/telemetry)}
% ================================================================================

\subsection{The Cross-Vendor Observability Problem}

When a saga spans five agents across three vendors, a 500\,ms latency spike is
impossible to diagnose from individual agent logs alone. The only solution is
distributed tracing: every message in the saga correlated under a single trace ID,
with each agent-to-agent hop recorded as a child span.

\texttt{@aroha-sdk/telemetry} implements the W3C Trace Context standard
(\texttt{traceparent}) as a first-class envelope field, propagated and signed
alongside every Aroha message.

\subsection{Trace Context Propagation}

The W3C \texttt{traceparent} format:
\begin{center}
  \texttt{00-\{traceId:32hex\}-\{spanId:16hex\}-\{flags:2hex\}}
\end{center}

The \texttt{traceparent} field is included in every \texttt{ArohaEnvelope} and
covered by the Ed25519 proof, so it cannot be tampered with or spoofed in transit.
Each receiving agent:
\begin{enumerate}[noitemsep]
  \item Extracts \texttt{traceparent} from the incoming envelope.
  \item Starts a child span (same \texttt{traceId}, new \texttt{spanId}).
  \item Propagates the same \texttt{traceId} on every outbound message for
        the duration of the saga.
\end{enumerate}

The result is a single distributed trace visible in Datadog, Grafana Tempo,
Jaeger, Honeycomb, or any OpenTelemetry-compatible backend --- with no custom
instrumentation beyond dropping the \texttt{createTracingMiddleware} into the
server's middleware chain.

\subsection{OpenTelemetry Integration}

When \texttt{@opentelemetry/api} is present in the host application,
\texttt{@aroha-sdk/telemetry} automatically propagates context into the OTel
\texttt{ContextAPI}, so Aroha spans appear as children of any existing application
traces. This requires zero additional configuration.

% ================================================================================
\section{Multi-Cloud Interoperability}
% ================================================================================

\subsection{The Vendor-Lock-In Gap}

Every major cloud vendor ships an agent protocol:

\begin{table}[ht]
  \centering
  \renewcommand{\arraystretch}{1.4}
  \begin{tabularx}{\textwidth}{@{} l l X @{}}
    \toprule
    \textbf{Vendor} & \textbf{Protocol} & \textbf{Structural limitation} \\
    \midrule
    Microsoft  & Azure Agent Protocol / A2A & Azure and M365 ecosystem only \\
    Google     & Agent-to-Agent (A2A)       & Vertex AI ecosystem only \\
    Amazon     & Bedrock Agents             & AWS ecosystem only \\
    Salesforce & AgentForce                 & Salesforce ecosystem only \\
    \bottomrule
  \end{tabularx}
  \caption{Major cloud agent protocols and their ecosystem scope}
\end{table}

These protocols are excellent within their ecosystems but share a structural
flaw: they are designed to maximise platform retention, not cross-vendor
interoperability. An Azure Copilot Studio agent has no standard way to discover
a Bedrock agent in a customer's AWS account, verify its identity without a
shared trust root, or execute an atomic saga that spans both.

\subsection{Aroha as the Inter-Cloud Trust Layer}

Aroha is positioned as the \textbf{inter-cloud trust spine} --- not a replacement
for any vendor's agent platform, but the layer that lets all of them talk to each
other.

\begin{itemize}[noitemsep]
  \item Each cloud's agent keeps its existing reasoning framework (LangChain,
        Semantic Kernel, Bedrock Converse API) unchanged.
  \item A Aroha bridge wrapper gives it a \texttt{did:aroha:} identity, a signed
        capability manifest, and the ability to participate in cross-vendor sagas.
  \item The \texttt{@aroha-sdk/a2a-bridge} handles Google A2A translation;
        \texttt{@aroha-sdk/mcp-bridge} handles Anthropic MCP;
        \texttt{@aroha-sdk/langchain-bridge} handles LangChain/Semantic Kernel agents.
\end{itemize}

\begin{table}[ht]
  \centering
  \renewcommand{\arraystretch}{1.4}
  \begin{tabularx}{\textwidth}{@{} l X @{}}
    \toprule
    \textbf{Claim} & \textbf{How Aroha delivers it} \\
    \midrule
    Vendor-neutral
      & Apache 2.0 license, W3C CG charter drafted (not yet proposed), no vendor controls the spec \\
    Cloud-agnostic identity
      & \texttt{did:aroha:} DIDs are not tied to any cloud IAM system \\
    No blockchain required
      & \texttt{InMemoryRegistry} + \texttt{PostgresSettlement} --- zero
        external dependencies \\
    Runs everywhere
      & Node.js 22+; deploys on Vercel, Cloudflare Workers, AWS Lambda,
        Azure Functions, GCP Cloud Run \\
    Existing agents unchanged
      & Bridge packages require zero modifications to the wrapped agent \\
    \bottomrule
  \end{tabularx}
  \caption{Aroha multi-cloud positioning}
\end{table}

% ================================================================================
\section{Post-Quantum Cryptography Roadmap}
% ================================================================================

\subsection{Design for No-Flag-Day Migration}

Aroha is designed for a zero-disruption migration to NIST post-quantum standards.
The entire signing surface is isolated to two source files:
\texttt{aroha-core/src/crypto/signing.ts} (Ed25519 sign/verify) and
\texttt{aroha-core/src/identity/did.ts} (keypair generation). All callers ---
\texttt{buildEnvelope}, \texttt{validateEnvelope}, credential issuance ---
are unchanged during migration.

\subsection{Algorithm Comparison}

\begin{table}[ht]
  \centering
  \renewcommand{\arraystretch}{1.4}
  \begin{tabularx}{\textwidth}{@{} l r r l @{}}
    \toprule
    \textbf{Algorithm} & \textbf{Public key} & \textbf{Signature} & \textbf{Security level} \\
    \midrule
    Ed25519 (current)       &     32 bytes &      64 bytes & 128-bit classical \\
    ML-DSA-44 (FIPS 204)    &  1,312 bytes &   2,420 bytes & NIST Level 2 (128-bit PQ) \\
    ML-DSA-65 (FIPS 204)    &  1,952 bytes &   3,293 bytes & NIST Level 3 (192-bit PQ) \\
    SLH-DSA-128s (FIPS 205) &     32 bytes &   7,856 bytes & NIST Level 1 (128-bit PQ, small key) \\
    \bottomrule
  \end{tabularx}
  \caption{Signing algorithm comparison. ML-DSA-44 is the recommended migration target.}
\end{table}

ML-DSA-44 (FIPS 204) is the recommended target: it provides NIST Level 2
post-quantum security at a key/signature size acceptable for HTTPS headers.

\subsection{Migration Procedure}

Migration proceeds in five steps with no mandatory flag day:

\begin{enumerate}[label=\textbf{Step \arabic*:}, leftmargin=*]
  \item \textbf{Dual key in DID Document} --- Add a second
        \texttt{verificationMethod} of type \texttt{Mldsa44VerificationKey2025}
        alongside the existing Ed25519 key. Legacy agents ignore unknown key types.

  \item \textbf{Dual-sign envelopes} --- During the transition window, sign each
        envelope with both keys. Place the Ed25519 proof in the \texttt{proof}
        field and the ML-DSA-44 proof in a \texttt{proofs} array. Legacy agents
        verify \texttt{proof} only; upgraded agents verify \texttt{proofs[0]}.

  \item \textbf{Swap signing module} --- Replace the Ed25519 import in
        \texttt{signing.ts} with \texttt{ml\_dsa44} from
        \texttt{@noble/post-quantum/ml-dsa}. No call-site changes required.

  \item \textbf{Transition window} --- Run dual-signing for a minimum of 7 days,
        or until all known peer agents confirm ML-DSA-44 support. New agents verify
        PQ signatures when present; legacy agents fall back to Ed25519.

  \item \textbf{Deprecate Ed25519} --- Remove the Ed25519 key from the DID
        Document, stop dual-signing, and remove Ed25519 verification from
        \texttt{ArohaServer}.
\end{enumerate}

% ================================================================================
\section{Conformance}
% ================================================================================

This section defines what constitutes a conformant Aroha implementation.
The key words \textbf{MUST}, \textbf{MUST NOT}, \textbf{SHOULD}, \textbf{MAY}
are used as defined in RFC 2119.

\begin{table}[ht]
  \centering
  \renewcommand{\arraystretch}{1.5}
  \begin{tabularx}{\textwidth}{@{} >{\bfseries}l X @{}}
    \toprule
    Requirement & Description \\
    \midrule
    MUST
      & Expose \texttt{GET /.well-known/aroha-agent.json} returning a valid
        DID Document signed by the agent's Ed25519 private key. \\
    MUST
      & Expose \texttt{POST /aroha/v1} accepting signed Aroha envelopes. \\
    MUST
      & Verify the Ed25519 signature on every inbound envelope before dispatching
        to application code. \\
    MUST
      & Validate \texttt{envelope.to} matches the agent's own DID. \\
    MUST
      & Validate \texttt{envelope.expires} and reject expired messages. \\
    MUST
      & Maintain a nonce store for the rolling 5-minute window and reject
        duplicate nonces. \\
    MUST
      & Implement \texttt{ArohaReserve}/\texttt{ArohaCommit}/\texttt{ArohaCancel}
        idempotently when acting as a saga participant. \\
    SHOULD
      & Use TLS 1.3 or higher on all connections. \\
    SHOULD
      & Register capability manifests with a Aroha-conformant registry. \\
    SHOULD
      & Sign capability manifests with the agent's Ed25519 private key. \\
    MAY
      & Support the \texttt{@aroha-sdk/credentials} RBAC extension. \\
    MAY
      & Support the \texttt{@aroha-sdk/settlement} payment extension. \\
    MAY
      & Support the \texttt{@aroha-sdk/commerce} commerce extension. \\
    MAY
      & Support streaming via \texttt{WebSocket} at \texttt{/aroha/stream}. \\
    MUST NOT
      & Define new top-level message types. Extension is done via the
        \texttt{body} field only. \\
    MUST NOT
      & Skip signature verification for any message, including messages from
        trusted-mesh peers (nonce and expiry checks are always enforced). \\
    \bottomrule
  \end{tabularx}
  \caption{Aroha conformance requirements}
\end{table}

% ================================================================================
\section{Governance}
% ================================================================================

\subsection{Specification Ownership}

Aroha is intended for a \textbf{W3C Community Group} as an open, royalty-free
specification; no Community Group has been proposed or joined yet, and nothing has been submitted. No single vendor owns the specification. Changes would require consensus
among Working Group participants as defined by the W3C Community Group Process.

\subsection{How to Contribute}

\begin{itemize}[noitemsep]
  \item \textbf{Specification issues}: Open a GitHub issue in the
        \texttt{aroha-spec} repository with the \texttt{spec} label.
  \item \textbf{New message types or layer changes}: Submit a GitHub Discussion
        tagged \texttt{proposal} before opening a pull request. Major changes require
        a two-week comment period.
  \item \textbf{Extension packages}: Extensions do not require Working Group
        approval. Publish to npm under any namespace; register with the community
        extension registry.
  \item \textbf{SDK contributions}: Standard open-source pull request process
        against the \texttt{aroha-sdk} monorepo. All PRs require passing CI and two
        reviewer approvals.
\end{itemize}

\subsection{Versioning Policy}

Aroha uses \textbf{semantic versioning} at the specification level:
\begin{itemize}[noitemsep]
  \item \textbf{Patch} versions (\texttt{1.0.x}): Non-normative clarifications,
        editorial corrections, errata.
  \item \textbf{Minor} versions (\texttt{1.x.0}): Backwards-compatible additions
        (new \textit{MAY} conformance items, new extension packages).
  \item \textbf{Major} versions (\texttt{x.0.0}): Backwards-incompatible changes
        to the kernel (message format, signing algorithm, DID method). Major versions
        require a 90-day deprecation window and dual-version support in all reference
        implementations.
\end{itemize}

% ================================================================================
\section{Glossary}
% ================================================================================

\begin{description}[labelwidth=5cm, leftmargin=!, font=\ttfamily]
  \item[Agent]           An autonomous software process that can send and receive
                         Aroha messages, expose capabilities, and participate in sagas.
  \item[CID]             Content Identifier. An IPFS self-describing hash used to
                         anchor capability manifests in a content-addressed store.
  \item[DID]             Decentralized Identifier. A W3C standard for
                         self-sovereign, cryptographically verifiable identifiers
                         (W3C DID Core 1.0, \texttt{did:aroha:} method).
  \item[DID Document]    The JSON-LD document at
                         \texttt{/.well-known/aroha-agent.json} containing an
                         agent's public key, service endpoint, and metadata.
  \item[EMA]             Exponential Moving Average. The formula used by the
                         \texttt{ArohaReputation} oracle:
                         $\text{score} = \lfloor\text{old} \times 0.85 +
                         \text{signal} \times 0.15\rfloor$.
  \item[EVM]             Ethereum Virtual Machine. The execution environment for
                         EVM-compatible smart contracts used by
                         \texttt{ArohaRegistry.sol} and \texttt{ArohaEscrow.sol}.
  \item[IPFS]            InterPlanetary File System. A content-addressed
                         distributed storage network used for capability manifest
                         hosting.
  \item[JSON-LD]         JSON Linked Data. A W3C standard for embedding semantic
                         context into JSON documents; used for Aroha envelopes and
                         DID Documents.
  \item[Orchestrator]    An agent that drives a saga: discovers providers,
                         issues Reserve/Commit/Cancel, and handles rollback.
  \item[Provider]        An agent that exposes capabilities in a signed manifest
                         and participates in sagas as a participant.
  \item[Saga]            A multi-step distributed transaction where each step is
                         reversible via a compensating action. Aroha implements the
                         saga pattern as a first-class protocol primitive.
  \item[VC]              Verifiable Credential. A W3C standard for tamper-evident
                         cryptographic credential exchange (W3C VC Data Model 2.0).
  \item[W3C CG]          W3C Community Group. An open-participation group that
                         develops specifications outside the formal W3C
                         Recommendation track. Aroha has not yet been submitted to one.
  \item[ZKP]             Zero-Knowledge Proof. A cryptographic technique allowing
                         one party to prove a statement is true without revealing
                         underlying data. Aroha VCs \textit{MAY} carry ZKP proofs
                         for privacy-preserving credential presentation.
\end{description}

% ================================================================================
\section{References}
% ================================================================================

\begin{enumerate}[label={[\arabic*]}, leftmargin=*, itemsep=4pt]
  \item W3C DID Working Group. \textit{Decentralized Identifiers (DIDs) v1.0}.
        W3C Recommendation. July 2022.
        \url{https://www.w3.org/TR/did-core/}

  \item W3C Verifiable Credentials Working Group. \textit{Verifiable Credentials
        Data Model v2.0}. W3C Candidate Recommendation. May 2024.
        \url{https://www.w3.org/TR/vc-data-model-2.0/}

  \item W3C JSON-LD Working Group. \textit{JSON-LD 1.1}.
        W3C Recommendation. July 2020.
        \url{https://www.w3.org/TR/json-ld11/}

  \item Bernstein, D. J. \textit{Ed25519: High-speed high-security signatures}.
        IACR ePrint 2011/368. Also specified as RFC 8032.
        \url{https://tools.ietf.org/html/rfc8032}

  \item NIST. \textit{FIPS 204 --- Module-Lattice-Based Digital Signature
        Standard (ML-DSA)}. August 2024.
        \url{https://csrc.nist.gov/pubs/fips/204/final}

  \item NIST. \textit{FIPS 205 --- Stateless Hash-Based Digital Signature
        Standard (SLH-DSA)}. August 2024.
        \url{https://csrc.nist.gov/pubs/fips/205/final}

  \item W3C. \textit{W3C Trace Context Level 1}.
        W3C Recommendation. February 2020.
        \url{https://www.w3.org/TR/trace-context/}

  \item Hohpe, G.; Woolf, B. \textit{Enterprise Integration Patterns}.
        Addison-Wesley, 2003. Saga pattern originally described by
        Garcia-Molina, H.; Salem, K. ``Sagas.'' ACM SIGMOD 1987.

  \item OASIS. \textit{PKCE --- Proof Key for Code Exchange by OAuth Public
        Clients}. RFC 7636. September 2015.
        \url{https://tools.ietf.org/html/rfc7636}

  \item Buterin, V. et al. \textit{Ethereum Yellow Paper}.
        \url{https://ethereum.github.io/yellowpaper/paper.pdf}

  \item Bradner, S. \textit{Key words for use in RFCs to Indicate Requirement
        Levels}. RFC 2119. March 1997.
        \url{https://tools.ietf.org/html/rfc2119}
\end{enumerate}

% ================================================================================
\section{Conclusion}
% ================================================================================

The autonomous agent economy is not a future possibility --- it is being assembled
right now, one proprietary API at a time. Without an open protocol, agent
interoperability will consolidate around whichever vendor reaches critical mass
first, creating a new form of platform lock-in at the infrastructure layer.

Aroha addresses this at the source. The six-layer stack provides a complete,
self-contained answer to the three infrastructure failures that make cross-vendor
agent systems unsafe today: agents can \textbf{discover} each other without
hardcoded URLs, \textbf{verify} each other without trusting platform-level claims,
and \textbf{transact} atomically without the risk of partial state or stranded funds.

The protocol is deliberately minimal at the kernel level. Layers 0--5 define
only what every conformant agent must implement: transport endpoints, a
\texttt{did:aroha} identity, a signed capability manifest, 15 message types, the
saga state machine, and preference context. Everything else --- payments,
permissions, commerce contracts, telemetry, bridge compatibility --- lives in
opt-in extension packages that the kernel never imports.

This design means Aroha is incrementally adoptable. An engineering team can ship
a \texttt{@aroha-sdk/core} provider in a single afternoon and connect it to the open
ecosystem immediately. They can add \texttt{@aroha-sdk/credentials} RBAC and
\texttt{@aroha-sdk/settlement} billing the following week. When commercial scale
demands it, they can graduate to on-chain escrow or Stripe settlement --- the
provider code does not change; only the injected strategy does.

The specification is built entirely on open standards: W3C Decentralised
Identifiers, W3C Verifiable Credentials, JSON-LD, Ed25519 cryptography, and
optionally EVM-compatible smart contracts. There are no proprietary formats,
no required SaaS registries, and no mandatory blockchain. Any operator can run
the full stack --- including a private HTTP registry, a private credential
registry, and the \texttt{@aroha-sdk/trusted-mesh} middleware for internal meshes ---
without any dependence on external infrastructure.

A W3C Community Group charter for Aroha is drafted but not yet proposed, and the
local runtime is not yet released. The reference TypeScript SDK covers all 12 packages
with complete documentation and over 100 automated tests. Multi-language SDKs
for Python, Go, and Java are available for teams that do not use TypeScript.

The window to establish open standards in a fast-moving ecosystem is always
short. Aroha is ready. The specification is open. The implementation is available.

\vspace{1.2cm}
\hrule
\vspace{0.5cm}
\noindent\small\textit{%
  Aroha Specification v1.0 --- Draft.
  Status: draft; not submitted to the W3C or any other standards body.
  Date: April 28, 2026.
  Authors: Aroha Working Group.
  License: Apache License 2.0 (see LICENSE and NOTICE in the repository).%
}

\end{document}